\begin{textsample}{NEG Dim 1 – vem\_ed\_24 – Score 15.00 – vem\_ed\_24/t129.txt}  \label{ex:f1_neg_135}
IVES Cesu 2024: Inteligência Artificial e inovação Gabriela Méndez, da Florida International University (EUA) Gabriela Méndez (Florida International University, EUA) exemplificou estratégias práticas para inovação em projetos COIL utilizando inteligência artificial.
Sua apresentação está disponível em http://go.fiu.edu/ives2024 A Inteligência Artificial Generativa funciona mediante o uso de um modelo de aprendizagem de máquina (machine learning).
Essa aprendizagem ocorre a partir de conteúdos criados por humanos (conjuntos de dados, padrões, relações).
A IA Generativa cria novos conteúdos a partir dos padrões aprendidos, é o caso dos chatbots (ChatGPT, Microsoft Copilot, Perplexity, Koala Chat e You.com), dos geradores de imagens (playground.com, Leonardo.ai, Stable Diffusion On-line), de apresentações (gamma.app), de vídeos (fliki.ai, genmo.ai) e de voz (elevenlabs.io).
Para interagir com a IA Generativa, é preciso redigir o prompt (pergunta, pedido, instrução) de maneira clara e específica, para obter respostas mais adequadas ao que estamos buscando.
É preciso incluir o contexto, \textbf{adicionar} detalhes, exemplos ou cenários relevantes para garantir que a mensagem se alinhe com seus objetivos. “Experimente e aperfeiçoe, pratique até que obtenha resultados com os quais esteja satisfeito”, recomenda Gabriela.
A IA Generativa é útil em projetos COIL, especialmente para brainstorming, design instrucional, comunicação e matchmaking entre professores de áreas que à primeira vista podem parecer distantes.
O Chat GPT pode ajudar a desenvolver ideias para projetos COIL, por exemplo. “Daqui a cinco anos a educação vai ser mais autônoma, mas os estudantes precisarão de alguém para \textbf{guiar} seus passos na instituição de ensino”, projeta Gabi, como é conhecida no mundo COIL. “Os alunos serão mais autodidatas, mas sempre com apoio do docente, que deve ter um papel de \textbf{coach}, de facilitador, de \textbf{guia}; para isso, temos que transformar nosso papel de educadores”.
Sobre o espírito questionador, imprescindível nas relações entre as inteligências humana e artificial, levantado na sessão de perguntas após a apresentação, Gabi respondeu: “Todos temos responsabilidade de tratar criticamente as ferramentas de IA”.
Se as referências da IA não representam os diferentes contextos sociais, é importante trazer visibilidade para esses contextos e, “para que isso aconteça, é preciso colaborar, utilizando essas ferramentas”.
Osvaldo Succi Junior finalizou as discussões sobre o tema com uma importante observação: “Especialmente nós que estamos no Sul Global temos que pensar muito cuidadosamente sobre essa questão da \textbf{representatividade}”.
Usar as ferramentas de IA com criticismo, sempre.

% matched lemmas: adicionar, coach, guia, guiar, representatividade
\end{textsample}
